%% kapa: exam=1 round=22 year=2011 period=2 form=A subject=accounting subject_ko=회계학 seq=12 number=12
%% source: Q-Net 감정평가사 기출문제 (cst00309) artlSeq=5036790 file=extracted/1008/5036790/2011년 제22회 감정평가사 1차시험 기출문제/2교시.hwp
%% conversion: kordoc 4.24.51 -> tools/gen1.py (auto); review=proofread-gpt-astra; flags=none
\begin{question}{12}
다음은 20×1년 말 (주)서울의 퇴직급여에 관한 자료이다.

\begin{framedtext}
\begin{tabular*}{\linewidth}{@{}l@{\extracolsep{\fill}}r@{}}
ㆍ확정급여채무의 현재가치 & ￦10,000,000 \\
ㆍ사외적립자산의 공정가치 & 9,000,000 \\
ㆍ미인식 보험수리적 이익 & 2,500,000 \\
ㆍ종업원의 예상평균잔여근무연수 & 10년
\end{tabular*}
\end{framedtext}

(주)서울이 보험수리적손익을 범위접근법방식에 따라 최소한만 인식한다고 할 때, 20×2년에 인식할 보험수리적이익은? \par
\begin{choices}
  \item ￦150,000
  \item ￦160,000
  \item ￦200,000
  \item ￦250,000
  \item ￦260,000
\end{choices}
\end{question}
\answer{1}
